\documentclass[10pt]{article}
\usepackage[a4paper, margin=0.72in]{geometry}
\usepackage{titlesec}
\usepackage{booktabs}
\usepackage{tabularx}
\usepackage{xcolor}
\usepackage{enumitem}
\usepackage{tgheros}
\usepackage[T1]{fontenc}
\usepackage{array}
\usepackage{textcomp}
\usepackage[hidelinks]{hyperref}
\usepackage{tikz}
\usetikzlibrary{shapes.rounded, arrows.meta, positioning, calc}
\renewcommand{\familydefault}{\sfdefault}

%% Colour palette
\definecolor{primary}{RGB}{15,52,96}        %% Deep petroleum navy
\definecolor{accent}{RGB}{198,124,0}        %% Oil amber
\definecolor{tier1}{RGB}{20,100,60}         %% Tier 1 green (start here)
\definecolor{tier2}{RGB}{140,80,0}          %% Tier 2 amber (build next)
\definecolor{tier3}{RGB}{80,80,80}          %% Tier 3 grey (longer term)
\definecolor{lightgrey}{gray}{0.85}
\definecolor{darkgrey}{gray}{0.38}

%% Section style
\titleformat{\section}{%
  \color{primary}\bfseries\large\raggedright
}{}{0em}{}[\color{lightgrey}{\titlerule[1.5pt]}\vspace{-5pt}]

\titleformat{\subsection}{%
  \color{accent}\bfseries\normalsize\raggedright
}{}{0em}{}

\setlength{\parindent}{0pt}
\setlength{\parskip}{4pt}
\hypersetup{colorlinks=true, urlcolor=primary}

%% =====================================================
\begin{document}

%% ----- Title block -----
\begin{center}
  \vspace*{0.15cm}
  {\color{primary}\LARGE\bfseries NEWCROSS AND PAN OCEAN GROUP}\\[5pt]
  {\color{primary}\large\bfseries AI and Data Roadmap --- Practical Intelligence
   for Nigerian Oil and Gas Operations}\\[10pt]
  {\color{darkgrey}\normalsize Version 1.0\quad\textbar\quad April 2026\quad\textbar\quad
   Companion to the Strategic Vision Paper}
\end{center}

\noindent\colorbox{primary!8}{%
\begin{minipage}{\dimexpr\textwidth-2\fboxsep}
  \vspace{5pt}
  \textbf{\color{primary}Guiding Principle.}\quad
  Don't build AI for AI's sake. Build AI for Nigerian realities. Every initiative
  in this document passes three tests before it is considered viable:
  \textbf{(1)} Does it solve a real, costly problem at the group today?
  \textbf{(2)} Can it be built with data and infrastructure that actually exists?
  \textbf{(3)} Could another Nigerian operator pay for it within three years?
  If any answer is no, the initiative is deprioritised.
  \vspace{5pt}
\end{minipage}}

\vspace{10pt}

%% ----------------------------------------------------------
\section{Nigerian Operational Reality --- Design Constraints}
%% ----------------------------------------------------------

{\small
Every AI system built for this group must be designed around these realities from
day one. They are not edge cases. They are the standard operating environment.
}

\vspace{4pt}
\small
\begin{tabularx}{\columnwidth}{@{}lX@{}}
\toprule
\textbf{Operational Reality} & \textbf{Design Implication} \\
\midrule
Intermittent internet at remote sites (Ogharefe, Owa Aladinma, AEPP corridor)
  & Offline-first; edge inference where possible; sync when connected \\[3pt]
Power instability across operational sites
  & Async processing; graceful degradation; no cloud-only dependency \\[3pt]
Low bandwidth at field locations
  & Lightweight models; local inference; minimal data transfer \\[3pt]
NUPRC/DPR regulatory framework including April 2026 GHG directive
  & Compliance logic baked into data models, not added as an afterthought \\[3pt]
PIA 2021 HCDT obligations (3\% of opex per entity)
  & HCDT tracking integrated into procurement and finance data pipelines
    from the first day of data collection \\[3pt]
Nigerian Content Act requirements
  & Spend and employment data structured for NCA reporting from day one \\[3pt]
Naira currency and FX exposure
  & Price and operate in Naira; minimise dollar-denominated cloud costs \\[3pt]
Data quality legacy issues across SAP and operational systems
  & Assume dirty data; build cleaning and validation pipelines before ML \\[3pt]
Trust and explainability culture in operations teams
  & AI recommendations must explain their reasoning to be accepted by
    engineers and plant managers \\[3pt]
Multi-entity group structure (three legal entities, partly shared operations)
  & All systems must support entity-level and consolidated views simultaneously \\
\bottomrule
\end{tabularx}
\normalsize

\vspace{10pt}

%% ----------------------------------------------------------
\section{Tier 1 --- Start Here: High ROI, Achievable with Existing Data}
%% ----------------------------------------------------------

\noindent\colorbox{tier1!8}{%
\begin{minipage}{\dimexpr\textwidth-2\fboxsep}
  \vspace{3pt}
  {\small\color{tier1}\textbf{Tier 1 Criteria:} Real problem with immediate cost impact.
  Buildable with data and infrastructure that exists today. Commercial potential clear.}
  \vspace{3pt}
\end{minipage}}

\vspace{6pt}

\subsection*{1.\ Predictive Maintenance --- OOGP Gas Plant (Highest Priority)}

{\small
\textbf{Problem.}
The Ovade-Ogharefe Gas Plant is the group's most complex and highest-consequence
asset. It is a large-scale cryogenic gas processing facility with 29 LPG storage
tanks. Equipment failures --- compressors, cryogenic units, storage pressure systems
--- are extremely costly and potentially dangerous. Maintenance is currently reactive:
the problem is discovered after it has already become a failure or near-failure.

\vspace{4pt}
\textbf{Solution.}
Instrument compressors, cryogenic processing units, and LPG storage tanks with
vibration, temperature, and pressure sensors. Train time-series anomaly detection
models on sensor readings combined with existing SAP maintenance logs. Output:
failure probability scores per asset class, recommended maintenance windows,
escalation alerts to the right personnel.

\vspace{4pt}
\textbf{Why now.}
The Phase 2 cryogenic expansion was completed in 2022. Maintenance pattern data
from the new cryogenic systems is still building. This is the ideal moment to
instrument fully and establish clean data before failure patterns are entrenched.
}

\vspace{4pt}
\small
\begin{tabularx}{\columnwidth}{@{}lX@{}}
\toprule
\textbf{Item} & \textbf{Detail} \\
\midrule
Data required     & SAP maintenance logs; SCADA feeds (existing); equipment OEM specs
                    and fault codes; sensor readings (new instrumentation required) \\
Tech stack        & Python, scikit-learn / PyTorch, TSAI or Prophet for time-series,
                    InfluxDB or TimescaleDB for sensor data, Azure IoT Hub for ingestion \\
Expected ROI      & Reduction in unplanned downtime (one avoided plant shutdown
                    justifies the full project); extended equipment lifespan; reduction
                    in emergency parts procurement \\
Commercial product & \textbf{FieldMaintain} SaaS --- every gas plant operator in Nigeria
                    has this problem. OOGP is the reference deployment. \\
\bottomrule
\end{tabularx}
\normalsize

\vspace{8pt}

\subsection*{2.\ Regulatory Reporting Automation --- NUPRC Returns and GHG Directive}

{\small
\textbf{Problem.}
Quarterly and annual NUPRC regulatory submissions require pulling data from multiple
systems (SAP, spreadsheets, manual records), formatting it to NUPRC templates, and
submitting under tight deadlines. This takes teams days per submission cycle across
OML-98, OML-147, and OOGP. The April 2026 NUPRC directive adds mandatory
measurement-based methane and GHG reporting templates --- making manual compliance
even more operationally demanding.

\vspace{4pt}
\textbf{Solution.}
Build data extraction pipelines from SAP and operational systems. Use structured
template-filling with LLM assistance to auto-populate required report sections,
including the new GHG/methane templates. Add validation rules against NUPRC submission
requirements. Output: draft submissions ready for human review and sign-off. For
GHG reporting, sensor data feeds directly into the reporting pipeline.
}

\vspace{4pt}
\small
\begin{tabularx}{\columnwidth}{@{}lX@{}}
\toprule
\textbf{Item} & \textbf{Detail} \\
\midrule
Data required     & SAP data; existing NUPRC templates; past approved submissions
                    (for validation); OOGP sensor data for GHG measurement \\
Tech stack        & Python, LangChain or direct LLM API (locally hosted for sensitive data),
                    SAP data connectors, PDF generation library \\
Expected ROI      & Days $\to$ hours per submission cycle; fewer errors; full audit trail;
                    GHG directive compliance without additional headcount \\
Commercial product & \textbf{RegulatoryReport} --- painful for every operator; foreign
                    vendors consistently fail on Nigerian regulatory specifics \\
\bottomrule
\end{tabularx}
\normalsize

\vspace{8pt}

\subsection*{3.\ Nigerian Content and HCDT Compliance Tracker}

{\small
\textbf{Problem.}
Tracking, calculating, and reporting Nigerian Content Act compliance (local spend
percentage, local employment, training documentation) is manual, error-prone, and
creates significant audit risk. Under PIA 2021, the HCDT obligation (3\% of opex)
adds a further layer of reporting across all three group entities --- each with
distinct cost bases and procurement records.

\vspace{4pt}
\textbf{Solution.}
Integrate with SAP procurement and HR data across all three entities. Auto-classify
suppliers and spend by Nigerian Content category using NLP classification models.
Calculate HCDT contributions per entity from opex data. Generate compliance reports
in NCA-required format with early-warning flags before audit cycles.
}

\vspace{4pt}
\small
\begin{tabularx}{\columnwidth}{@{}lX@{}}
\toprule
\textbf{Item} & \textbf{Detail} \\
\midrule
Data required     & SAP procurement records; HR employee and training data; vendor
                    master lists; opex ledger per entity \\
Tech stack        & Python, spaCy or fine-tuned BERT for text classification,
                    SAP data connectors, report generation pipeline \\
Expected ROI      & Audit readiness; reduced compliance risk; staff time savings;
                    HCDT obligations tracked automatically \\
Commercial product & \textbf{NigerianContent.ai} --- mandatory for all operators;
                    PIA 2021 raised the stakes significantly \\
\bottomrule
\end{tabularx}
\normalsize

\vspace{8pt}

\subsection*{4.\ AEPP Integrity and Third-Party Tariff Reconciliation}

{\small
\textbf{Problem.}
The Amukpe-Escravos Pipeline serves third-party injectors on a tariff basis ---
meaning metering accuracy and theft detection are direct revenue issues, not merely
operational ones. 67 km through the Niger Delta with CHDD construction still presents
exposure. Tariff reconciliation for third-party injectors is currently paper-intensive
and manually reconciled.

\vspace{4pt}
\textbf{Solution.}
Deploy flow meter anomaly detection across all injection and delivery points.
Satellite imagery change detection (monthly comparison) along the 67 km Right of Way
to identify new construction, dig activity, or vegetation clearance near the pipeline.
Automated metering and tariff billing for injector volumes with full audit trail.
GPS-coordinated alert system for field response teams.
}

\vspace{4pt}
\small
\begin{tabularx}{\columnwidth}{@{}lX@{}}
\toprule
\textbf{Item} & \textbf{Detail} \\
\midrule
Data required     & SCADA flow meter readings; AEPP GPS coordinates; satellite imagery
                    (Planet Labs, Airbus Defence, or NASRDA preferred as Nigerian-first);
                    existing tariff and injection records from SAP \\
Tech stack        & Python, GDAL/rasterio for satellite imagery, statistical process
                    control for flow anomaly detection, GIS integration \\
Expected ROI      & Crude theft prevention; metering accuracy directly protects tariff
                    revenue; removes manual reconciliation overhead \\
Commercial product & \textbf{FieldWatch} --- every operator with long-haul Niger Delta
                    pipelines faces identical problems; AEPP is a 67 km reference site \\
\bottomrule
\end{tabularx}
\normalsize

\vspace{8pt}

\subsection*{5.\ HSE Digital Permit-to-Work and Incident Reporting}

{\small
\textbf{Problem.}
OOGP operates cryogenic processing equipment and 29 LPG storage tanks in a
high-consequence environment. Permit-to-work (PTW) systems are paper-based.
Incident reporting and near-miss tracking are manual. NUPRC has increasing
expectations around HSE reporting standards for high-consequence facilities.

\vspace{4pt}
\textbf{Solution.}
Replace paper PTW with a digital system configured for the OOGP environment.
Integrate incident reporting and near-miss tracking with a searchable database.
Generate HSE reports automatically for NUPRC submissions. Mobile-accessible for
field personnel; offline-capable for areas without reliable connectivity.
}

\vspace{4pt}
\small
\begin{tabularx}{\columnwidth}{@{}lX@{}}
\toprule
\textbf{Item} & \textbf{Detail} \\
\midrule
Data required     & Existing PTW records (digitise backlog for training); incident
                    and near-miss historical records; NUPRC HSE reporting templates \\
Tech stack        & React Native (mobile + offline-first), FastAPI backend,
                    PostgreSQL, Azure sync \\
Expected ROI      & Liability reduction; NUPRC compliance readiness; searchable
                    incident history for trend analysis \\
Commercial product & \textbf{HSETrack} --- foreign HSE software does not understand
                    Nigerian site conditions or NUPRC requirements \\
\bottomrule
\end{tabularx}
\normalsize

\vspace{10pt}

%% ----------------------------------------------------------
\section{Tier 2 --- Build Next: Medium Complexity, Requires Data Preparation}
%% ----------------------------------------------------------

\noindent\colorbox{tier2!8}{%
\begin{minipage}{\dimexpr\textwidth-2\fboxsep}
  \vspace{3pt}
  {\small\color{tier2}\textbf{Tier 2 Criteria:}
  High strategic value but requires data pipeline work from Tier 1 to be in place first.
  Begin planning in parallel; begin building once data foundations are established.}
  \vspace{3pt}
\end{minipage}}

\vspace{6pt}

\subsection*{6.\ Document Intelligence --- Well Logs, Drilling Reports, and Contracts}

{\small
\textbf{Problem.}
OML-98 has been producing since 1976. That is fifty years of drilling reports, well logs,
geological surveys, production records, and contracts --- most of which exist as PDFs,
scanned paper, or disconnected digital files. Finding specific data (a well's pressure
history, a vendor's contract terms, a drilling decision rationale) takes hours or days.
The group's institutional memory is largely inaccessible.

\vspace{4pt}
\textbf{Solution.}
OCR pipeline to digitise historical paper documents. LLM-based document classification
and metadata extraction. Semantic search (RAG architecture) so engineers can query
natural language questions against the full document archive. Unlock fifty years of
operational knowledge.
}

\vspace{4pt}
\small
\begin{tabularx}{\columnwidth}{@{}lX@{}}
\toprule
\textbf{Item} & \textbf{Detail} \\
\midrule
Data required     & Existing document archives (digital and physical). Digitisation
                    of paper records is the primary prerequisite. \\
Tech stack        & Python, Tesseract / AWS Textract for OCR, LangChain + vector
                    database (pgvector or Weaviate), locally-hosted LLM (Llama3 or
                    Mistral) for sensitive data \\
Security note     & Well data and production history is commercially sensitive.
                    Use locally-hosted LLMs only for this application. \\
Expected ROI      & Engineer productivity; faster decision-making; institutional
                    knowledge preserved beyond individual tenure \\
\bottomrule
\end{tabularx}
\normalsize

\vspace{8pt}

\subsection*{7.\ Unified Operations Dashboard with Anomaly Detection}

{\small
\textbf{Problem.}
Management has no real-time view across OML-98, OML-147, OOGP, and AEPP operations
simultaneously. Reports are generated periodically, so problems surface late. The
three-entity group structure means consolidated visibility is even harder to achieve
manually.

\vspace{4pt}
\textbf{Solution.}
Unified data integration layer pulling from SAP, SCADA, and operational systems across
all entities. Real-time dashboard (web and mobile-accessible) with configurable KPIs
per entity and consolidated group view. Anomaly detection layer flagging deviations
from expected parameters. Alert system with escalation logic.
}

\vspace{4pt}
\small
\begin{tabularx}{\columnwidth}{@{}lX@{}}
\toprule
\textbf{Item} & \textbf{Detail} \\
\midrule
Data required     & SCADA feeds; SAP operational data; logistics and dispatch records;
                    gas plant throughput data; AEPP metering records \\
Tech stack        & Python (FastAPI backend), React or React Native (frontend),
                    TimescaleDB, statistical anomaly detection layer \\
Expected ROI      & Faster management response to operational issues; group-level
                    visibility that currently does not exist in any form \\
Commercial product & \textbf{WellWatch} SaaS --- configurability for multi-entity,
                    multi-block operations is the key differentiator \\
\bottomrule
\end{tabularx}
\normalsize

\vspace{8pt}

\subsection*{8.\ Oil Marketing Intelligence}

{\small
\textbf{Problem.}
Crude lifting schedules, buyer allocation, and NNPC PSC equity crude reconciliation
on OML-147 are managed manually. The domestic crude market has materially changed with
Dangote refinery's entry. Oil marketing teams are making significant commercial
decisions without a systematic data environment to support them.

\vspace{4pt}
\textbf{Solution.}
Lifting schedule management and buyer allocation tracking system. PSC equity crude
reconciliation with NNPC, integrating with the OML-147 PSC reporting pipeline.
Domestic and export pricing dashboard tracking refinery netbacks, regional differentials,
and historical buyer performance.
}

\vspace{4pt}
\small
\begin{tabularx}{\columnwidth}{@{}lX@{}}
\toprule
\textbf{Item} & \textbf{Detail} \\
\midrule
Data required     & Lifting records; buyer contracts; NNPC PSC equity crude schedules;
                    Platts or Argus price feeds (API) \\
Tech stack        & Python backend, React frontend, PostgreSQL, price feed integrations \\
Expected ROI      & Lifting schedule accuracy; reduced reconciliation disputes;
                    better-informed commercial decisions in a shifting market \\
\bottomrule
\end{tabularx}
\normalsize

\vspace{10pt}

%% ----------------------------------------------------------
\section{Tier 3 --- Longer-Term Vision: Requires Significant Data Foundation}
%% ----------------------------------------------------------

\noindent\colorbox{tier3!10}{%
\begin{minipage}{\dimexpr\textwidth-2\fboxsep}
  \vspace{3pt}
  {\small\color{tier3}\textbf{Tier 3 Criteria:}
  High long-term value but requires 12--18 months of upstream data pipeline work
  before ML is viable. Plan architecture now; do not attempt to build now.}
  \vspace{3pt}
\end{minipage}}

\vspace{6pt}

\subsection*{9.\ Production Optimisation --- OML-98 and OML-147}

{\small
\textbf{Problem.}
Well and field performance across both OML blocks is managed based on engineering
experience and periodic review rather than continuous optimisation. Small improvements
in choke settings, injection schedules, and pressure management across two producing
blocks translate to meaningful revenue.

\vspace{4pt}
\textbf{Solution.}
Optimisation models trained on historical production data. Recommendations for well
scheduling, injection rates, and choke settings. Simulation environment for testing
operational decisions without field risk.

\vspace{4pt}
\textbf{Prerequisite.}
Requires significant data quality improvement work from Tier 1 and Tier 2 initiatives.
This is an 18-month build, not a starting point.
}

\vspace{8pt}

\subsection*{10.\ Procurement and Supply Chain Optimisation}

{\small
\textbf{Problem.}
Procurement for remote field operations involves significant lead times, FX exposure
on imported parts, and frequent emergency orders at premium cost. Vendor performance
is not systematically tracked. Three entities with partly overlapping procurement
needs creates further inefficiency.

\vspace{4pt}
\textbf{Solution.}
Demand forecasting for spare parts based on maintenance schedules and equipment age
(feeding directly from the Tier 1 predictive maintenance system). Vendor risk scoring
from historical delivery performance. FX exposure dashboard for imported goods.
Consolidated group procurement view to identify cross-entity opportunities.
}

\vspace{10pt}

%% ----------------------------------------------------------
\section{AI Design Principles for Group Deployments}
%% ----------------------------------------------------------

\small
\begin{tabularx}{\columnwidth}{@{}lX@{}}
\toprule
\textbf{Principle} & \textbf{Rationale and Application} \\
\midrule
Offline-first
  & Core AI functions must produce useful output without reliable internet.
    OOGP, Ogharefe, and the AEPP corridor are not reliably connected environments. \\[3pt]
Edge inference where possible
  & Run model inference locally at the asset; send only anomaly alerts and summaries
    to the cloud. Reduces bandwidth, latency, and cloud cost. \\[3pt]
Explainable outputs
  & Every AI recommendation must include a plain-language explanation.
    ``Compressor Unit 3: 78\% failure probability within 14 days based on bearing
    vibration trend and last maintenance date.'' Not just a score. \\[3pt]
Locally hosted for sensitive data
  & Well data, production history, commercial contracts, and lifting records must
    not leave the group's Azure tenant or on-premise servers. Use locally-hosted
    LLMs (Llama3, Mistral) for these applications. \\[3pt]
Assume dirty data
  & Every pipeline includes a validation and cleaning layer before data reaches
    any model. Data quality is logged and reported; not silently fixed. \\[3pt]
Multi-entity from day one
  & Every model and dashboard is built to operate at both entity and consolidated
    group level. Single-entity models will never be commercially reusable. \\[3pt]
Compliance logic is first-class
  & NUPRC reporting requirements, HCDT calculations, and NCA classifications are
    embedded in data models, not bolted on as reporting features later. \\
\bottomrule
\end{tabularx}
\normalsize

\vspace{10pt}

%% ----------------------------------------------------------
\section{Implementation Timeline}
%% ----------------------------------------------------------

\begin{center}
\begin{tikzpicture}[
  node distance=0.3cm,
  phase/.style={draw, rounded corners=5pt, minimum width=15cm,
                minimum height=1.0cm, align=left, font=\small, inner sep=8pt},
  lbl/.style={font=\small\bfseries, color=primary},
]

\node[phase, fill=tier1!10, draw=tier1!60] (y1q1)
  {\textbf{Year 1 | Q1--Q2}\quad
   Data Audit (SAP, SCADA, OOGP sensors, AEPP metering records) $\cdot$
   OOGP Predictive Maintenance Pilot (compressors, one asset class) $\cdot$
   Regulatory Reporting Automation MVP (OML-98 + OML-147 NUPRC returns + GHG templates) $\cdot$
   AEPP Tariff Reconciliation System $\cdot$ HSE Digital PTW (OOGP)};

\node[phase, fill=tier1!5, draw=tier1!40, below=of y1q1] (y1q3)
  {\textbf{Year 1 | Q3--Q4}\quad
   NCA + HCDT Tracker Phase 1 $\cdot$
   Unified Operations Dashboard (OML-98 + OML-147 + OOGP + AEPP) $\cdot$
   Predictive Maintenance Expansion (full OOGP cryogenic units + LPG storage) $\cdot$
   GHG Monitoring Pipeline (sensor-to-report, NUPRC templates)};

\node[phase, fill=tier2!8, draw=tier2!50, below=of y1q3] (y2q1)
  {\textbf{Year 2 | Q1--Q2}\quad
   Document Intelligence (OML-98 well logs from 1976 --- digitise 50 years) $\cdot$
   AEPP Pipeline Security Pilot (flow anomaly + satellite encroachment) $\cdot$
   Oil Marketing Intelligence (lifting + PSC equity reconciliation) $\cdot$
   Package \textbf{RegulatoryReport} for first external customer};

\node[phase, fill=tier2!5, draw=tier2!30, below=of y2q1] (y2q3)
  {\textbf{Year 2 | Q3--Q4}\quad
   \textbf{WellWatch} dashboard commercialisation $\cdot$
   \textbf{NigerianContent.ai} launch to first external operator $\cdot$
   OML-147 production optimisation data collection complete (ML modelling begins) $\cdot$
   API-first architecture and multi-tenant infrastructure established};

\node[phase, fill=tier3!8, draw=tier3!40, below=of y2q3] (y3)
  {\textbf{Year 3}\quad
   Production Optimisation (OML-98 + OML-147 ML models; simulation environment) $\cdot$
   Supply Chain Optimisation (demand forecasting; vendor risk scoring) $\cdot$
   Full product commercialisation (\textbf{FieldMaintain}, \textbf{FieldWatch},
   \textbf{HSETrack}) $\cdot$ Group data lake fully operational};

\end{tikzpicture}
\end{center}

\vspace{10pt}

%% ----------------------------------------------------------
\section{Data and Infrastructure Prerequisites}
%% ----------------------------------------------------------

{\small
No AI initiative succeeds without the data foundation beneath it.
The following are non-negotiable prerequisites before any machine learning
project begins.
}

\vspace{4pt}
\noindent\colorbox{primary!8}{%
\begin{minipage}{\dimexpr\textwidth-2\fboxsep}
  \vspace{5pt}
  \textbf{\color{primary}Communications Infrastructure as a Transformation Dependency.}\quad
  The group's existing LAN, WAN, voice, and remote rig data communications
  infrastructure --- spanning Lagos HQ, Ogharefe, Owa Aladinma, and Newcross
  E\&P rig sites --- is the physical foundation for every data initiative in
  this roadmap. A degraded WAN link to Ogharefe does not just reduce
  productivity: it silences SCADA feeds and creates blind spots in operational
  monitoring. A failed VSAT link to a remote rig means drilling parameter data
  stops transmitting and decisions are made on stale information. Communications
  infrastructure health is therefore a dependency of the transformation
  programme, not a separate IT concern. A connectivity audit across all sites
  must run in parallel with the data audit in Month 1.
  \vspace{5pt}
\end{minipage}}

\vspace{6pt}
\small
\begin{tabularx}{\columnwidth}{@{}lXl@{}}
\toprule
\textbf{Prerequisite} & \textbf{What It Means} & \textbf{Timeline} \\
\midrule
Communications and connectivity audit
  & Assess WAN, LAN, VSAT, and rig data link capacity, reliability, and
    redundancy across all sites (Lagos HQ, Ogharefe, Owa Aladinma, rig
    sites). Identify gaps that would prevent reliable data pipeline operation.
    Runs in parallel with the data audit.
  & Month 1--2 \\[3pt]
Data audit
  & Assess quality, completeness, and accessibility of SAP, SCADA, and
    operational data across all three entities
  & Month 1--2 \\[3pt]
Data governance framework
  & Define ownership, quality standards, access controls, and retention
    policies across the group; NDPR compliance built in
  & Month 1--3 \\[3pt]
Sensor instrumentation --- OOGP
  & Install vibration, temperature, and pressure sensors on priority
    compressor and cryogenic units where not already in place
  & Month 1--4 \\[3pt]
SAP data connectors
  & Build or procure reliable, authenticated API connections to SAP
    financial, maintenance, procurement, and HR modules
  & Month 1--3 \\[3pt]
Azure IoT Hub configuration
  & Set up IoT Hub for OOGP sensor ingestion; configure time-series
    database (InfluxDB or TimescaleDB) for sensor data storage
  & Month 2--4 \\[3pt]
Satellite imagery subscription
  & NASRDA (Nigerian-first), Planet Labs, or Airbus Defence for
    AEPP corridor change detection
  & Month 3--6 \\
\bottomrule
\end{tabularx}
\normalsize

\vspace{10pt}

%% ----- Footer -----
\vspace{6pt}
{\small\color{darkgrey}
\textbf{Newcross and Pan Ocean Group of Companies}\enspace|\enspace
AI and Data Roadmap --- Version 1.0\enspace|\enspace April 2026\\
Companion document: \textit{Digital Transformation Strategic Vision ---
Newcross and Pan Ocean Group (April 2026)}.\\
All AI use cases are grounded in available data, existing infrastructure,
and Nigerian operational reality --- not research projects.
}

\end{document}
